\section{결론}

본 논문은 멀티-툴 도구 선택에서 실패하는 것은 steering 자체가 아니라 \emph{정지 K-style steering의 배치}라는 점을 보였다. 핵심은 K를 더 세게 쓰는 것이 아니라, K를 초기 레이어의 정보 각인에만 사용하고 전 레이어의 coverage 조정은 Q에 맡기는 것이다. 이 설계는 MetaTool Subtask4에서 정지 K의 $-4.57$pp 붕괴를 같은 기저에서 $+2.08$pp 회복으로 뒤집고, $\tau^2$ 세 도메인에서 $+5.98$/$+26.76$/$+3.84$pp를 기록한다. retail을 horizon으로 분해하면 layer-adaptive와 Q-only가 short vs long horizon에서 다른 우세 구간을 가지므로, 두 연산자는 경쟁 관계라기보다 같은 family 안의 regime-specific 끝점으로 읽는 것이 맞다. 정리~\ref{thm:no-memory}과 따름정리~\ref{cor:cache-divergence}는 왜 정지 K가 멀티-툴에서 반복 오류를 피하기 어렵고, 왜 per-step에는 동등한 Q가 multi-step에서는 더 안정적일 수 있는지를 설명한다. 남은 과제는 분명하다. protocol-matched SEKA/AdaSEKA 비교로 prior-work closure를 끝내고, PCA-of-K와 format-controlled Llama audit으로 ontology specificity와 failure mechanism을 더 직접적으로 닫아야 한다. 그럼에도 현재 결과가 이미 말해 주는 바는 충분히 선명하다. 멀티-툴에서 중요한 것은 ``더 강한 정지 steering''이 아니라 \emph{어느 깊이에서 어떤 쪽을 조작할 것인가}다.
